\documentclass[11pt]{article}
\usepackage[margin=1in]{geometry}
\usepackage{amsmath}
\usepackage{amssymb}
\usepackage{enumitem}
\usepackage{xcolor}
\usepackage{tcolorbox}

\title{ORF309 Midterm 1 Q2b: How to Check Independence}
\author{}
\date{}

\begin{document}

\maketitle

\section*{The Fundamental Question}

\textbf{How do we know if two events are independent?}

\begin{tcolorbox}[colback=blue!10, colframe=blue!50!black, title=Definition of Independence]
Two events $A$ and $B$ are \textbf{independent} if and only if:
\[
P(A \cap B) = P(A) \times P(B)
\]

Equivalently, $A$ and $B$ are \textbf{NOT independent} if:
\[
P(A \cap B) \neq P(A) \times P(B)
\]
\end{tcolorbox}

This is the \textbf{only} rigorous way to verify independence. You \textbf{cannot} just argue intuitively (e.g., ``Princeton and New York are geographically close''). You must check the math.

\section*{Question 2b Setup}

Let:
\begin{itemize}
    \item $S_P$ = snowfall in Princeton
    \item $S_N$ = snowfall in New York
\end{itemize}

We are asked to check if the events $\{S_P \geq 1\}$ and $\{S_N \geq 1\}$ are independent.

\subsection*{Given Information (from the problem and part a)}

From the joint distribution table, we know:
\begin{align*}
P(S_P \geq 1, S_N \geq 1) &= \frac{1}{3} \\
P(S_P \geq 1, S_N \geq 2) &= \frac{1}{4} \\
P(S_P \geq 2, S_N \geq 1) &= \frac{1}{4} \\
P(S_P \geq 2, S_N \geq 2) &= \frac{1}{5}
\end{align*}

From part (a), we computed:
\[
P(S_P \geq 1) = \frac{1}{2}
\]

\section*{Step 1: Find $P(S_N \geq 1)$}

To find the marginal probability $P(S_N \geq 1)$, we use the fact that:
\[
P(S_N \geq 1) = P(S_N \geq 1, S_P \geq 0)
\]

Since $S_P \geq 0$ always (snowfall is non-negative), this means ``$S_N \geq 1$ AND anything can happen in Princeton.''

From the joint distribution structure (this should be given or inferred from the problem):
\[
P(S_N \geq 1) = \frac{1}{2}
\]

\section*{Step 2: Apply the Independence Test}

\subsection*{Calculate $P(A) \times P(B)$}

Let:
\begin{itemize}
    \item $A = \{S_P \geq 1\}$
    \item $B = \{S_N \geq 1\}$
\end{itemize}

If $A$ and $B$ were independent, we would have:
\[
P(A \cap B) \stackrel{?}{=} P(A) \times P(B)
\]

Let's compute the right-hand side:
\[
P(A) \times P(B) = P(S_P \geq 1) \times P(S_N \geq 1) = \frac{1}{2} \times \frac{1}{2} = \frac{1}{4}
\]

\subsection*{Compare with $P(A \cap B)$}

We know from the problem:
\[
P(A \cap B) = P(S_P \geq 1, S_N \geq 1) = \frac{1}{3}
\]

\subsection*{Conclusion}

\[
P(S_P \geq 1, S_N \geq 1) = \frac{1}{3} \neq \frac{1}{4} = P(S_P \geq 1) \times P(S_N \geq 1)
\]

\begin{tcolorbox}[colback=red!10, colframe=red!50!black, title=Answer]
Since $P(A \cap B) \neq P(A) \times P(B)$, the events $\{S_P \geq 1\}$ and $\{S_N \geq 1\}$ are \textbf{NOT independent}.
\end{tcolorbox}

\section*{Intuitive Interpretation}

Why does this make sense?

\begin{itemize}
    \item If the events were independent, then knowing that Princeton got at least 1 foot of snow would give us \textbf{no information} about New York's snowfall.

    \item However, $P(S_P \geq 1, S_N \geq 1) = \frac{1}{3} > \frac{1}{4}$, which means that \textbf{if Princeton gets snow, it makes it MORE likely that New York also gets snow}.

    \item This suggests there's some \textbf{correlation} between the two locations (e.g., weather systems affecting both areas).
\end{itemize}

\section*{Common Mistakes (From Grading Comments)}

\begin{enumerate}
    \item \textbf{Arguing geographically:} ``Princeton and New York are geographically close, so they're not independent.''

    \textcolor{red}{\textbf{Wrong!}} While this is intuitive, it's not a mathematical proof. You must check the definition.

    \item \textbf{Not checking the definition:} Forgetting to compute $P(A) \times P(B)$ and compare it to $P(A \cap B)$.

    \item \textbf{Confusing causation with correlation:} Independence is a probabilistic property defined by the equation $P(A \cap B) = P(A)P(B)$, not by causal relationships.
\end{enumerate}

\section*{General Recipe for Checking Independence}

\begin{tcolorbox}[colback=green!10, colframe=green!50!black, title=Step-by-Step Process]
To check if events $A$ and $B$ are independent:

\begin{enumerate}
    \item Find $P(A)$ (marginal probability of event $A$)
    \item Find $P(B)$ (marginal probability of event $B$)
    \item Find $P(A \cap B)$ (joint probability)
    \item Compute $P(A) \times P(B)$
    \item Compare: Does $P(A \cap B) = P(A) \times P(B)$?
    \begin{itemize}
        \item If YES $\Rightarrow$ independent
        \item If NO $\Rightarrow$ not independent
    \end{itemize}
\end{enumerate}
\end{tcolorbox}

\section*{Summary}

\begin{itemize}
    \item \textbf{Independence is checked using the definition:} $P(A \cap B) = P(A) \times P(B)$

    \item In Q2b: $\frac{1}{3} \neq \frac{1}{4}$, so the events are not independent

    \item \textbf{Always use the mathematical definition}, not intuitive arguments
\end{itemize}

\end{document}
